\textbf{Phần 1. Trắc nghiệm lựa chọn}
\setcounterpageref{ex}{0}

\begin{ex} %[3D2V70]
Tính giá trị của phép tính: $18\,240 \times 4$.
\choice
{$72\,860$}
{$73\,960$}
{$71\,960$}
{\True $72\,960$}
\loigiai{
$18\,240 \times 4 = 72\,960$.
\textbf{Đáp án của câu hỏi là D}
}
\end{ex}

\begin{ex} %[3D2V70]
Một phân xưởng mỗi ngày may được $12\,500$ chiếc khẩu trang. Hỏi trong 6 ngày phân xưởng đó may được bao nhiêu chiếc khẩu trang?
\choice
{$70\,000$ chiếc}
{\True $75\,000$ chiếc}
{$80\,000$ chiếc}
{$72\,500$ chiếc}
\loigiai{
$12\,500 \times 6 = 75\,000$ chiếc.
\textbf{Đáp án của câu hỏi là B}
}
\end{ex}

\begin{ex} %[3D2V70]
Tìm số tự nhiên $x$, biết: $x : 5 = 16\,400$.
\choice
{$80\,000$}
{$84\,000$}
{\True $82\,000$}
{$81\,000$}
\loigiai{
$x = 16\,400 \times 5 = 82\,000$.
\textbf{Đáp án của câu hỏi là C}
}
\end{ex}

\begin{ex} %[3D2V70]
Một hình vuông có cạnh dài $14\,250\text{ m}$. Chu vi hình vuông đó là
\choice
{\True $57\,000\text{ m}$}
{$56\,000\text{ m}$}
{$58\,000\text{ m}$}
{$57\,500\text{ m}$}
\loigiai{
Chu vi: $14\,250 \times 4 = 57\,000\text{ (m)}$.
\textbf{Đáp án của câu hỏi là A}
}
\end{ex}

\begin{ex} %[3D2V70]
Tính nhanh giá trị: $12\,500 \times 2 + 12\,500 \times 6$.
\choice
{$90\,000$}
{$80\,000$}
{$125\,000$}
{\True $100\,000$}
\loigiai{
$12\,500 \times (2 + 6) = 12\,500 \times 8 = 100\,000$.
\textbf{Đáp án của câu hỏi là D}
}
\end{ex}

\begin{ex} %[3D2V70]
Một đoàn gồm 3 xe tải, mỗi xe chở $24\,500\text{ kg}$ xi măng. Người ta đã dỡ xuống $35\,000\text{ kg}$. Số xi măng còn lại trên xe là
\choice
{$40\,500\text{ kg}$}
{\True $38\,500\text{ kg}$}
{$37\,500\text{ kg}$}
{$39\,000\text{ kg}$}
\loigiai{
Tổng xi măng: $24\,500 \times 3 = 73\,500\text{ kg}$. Còn lại: $73\,500 - 35\,000 = 38\,500\text{ (kg)}$.
\textbf{Đáp án của câu hỏi là B}
}
\end{ex}

\begin{ex} %[3D2V70]
Tính giá trị biểu thức: $A = 21\,340 \times 3 - 18\,500$.
\choice
{$44\,520$}
{$46\,520$}
{\True $45\,520$}
{$45\,020$}
\loigiai{
$21\,340 \times 3 = 64\,020$; $64\,020 - 18\,500 = 45\,520$.
\textbf{Đáp án của câu hỏi là C}
}
\end{ex}

\begin{ex} %[3D2V70]
Một công ty nhập về 4 lô hàng, mỗi lô có $15\,200$ hộp sữa. Sau khi xuất bán $28\,000$ hộp, công ty còn lại bao nhiêu hộp sữa?
\choice
{\True $32\,800$ hộp}
{$30\,800$ hộp}
{$34\,800$ hộp}
{$33\,000$ hộp}
\loigiai{
Tổng sữa: $15\,200 \times 4 = 60\,800$ hộp. Còn lại: $60\,800 - 28\,000 = 32\,800$ hộp.
\textbf{Đáp án của câu hỏi là A}
}
\end{ex}

\begin{ex} %[3D2V70]
Nếu tăng một thừa số thêm 3 đơn vị và giữ nguyên thừa số kia bằng $25\,000$ thì tích tăng thêm bao nhiêu đơn vị?
\choice
{$50\,000$}
{$25\,000$}
{$100\,000$}
{\True $75\,000$}
\loigiai{
Tích tăng thêm: $25\,000 \times 3 = 75\,000$ đơn vị.
\textbf{Đáp án của câu hỏi là D}
}
\end{ex}

\begin{ex} %[3D2V70]
Tính: $11\,111 \times 9$. Kết quả là
\choice
{$99\,990$}
{\True $99\,999$}
{$100\,000$}
{$98\,999$}
\loigiai{
$11\,111 \times 9 = 99\,999$.
\textbf{Đáp án của câu hỏi là B}
}
\end{ex}

\begin{ex} %[3D2V70]
Một huyện có 5 xã vùng cao, mỗi xã được cấp $18\,500\text{ kg}$ gạo hỗ trợ mùa giáp hạt. Tổng số gạo đã cấp cho cả 5 xã là
\choice
{$90\,000\text{ kg}$}
{$91\,500\text{ kg}$}
{\True $92\,500\text{ kg}$}
{$93\,500\text{ kg}$}
\loigiai{
$18\,500 \times 5 = 92\,500\text{ (kg)}$.
\textbf{Đáp án của câu hỏi là C}
}
\end{ex}

\begin{ex} %[3D2V70]
Một cuộn dây thép dài $32\,400\text{ m}$. Ba cuộn dây thép như thế dài bao nhiêu mét?
\choice
{\True $97\,200\text{ m}$}
{$96\,200\text{ m}$}
{$98\,200\text{ m}$}
{$95\,200\text{ m}$}
\loigiai{
$32\,400 \times 3 = 97\,200\text{ (m)}$.
\textbf{Đáp án của câu hỏi là A}
}
\end{ex}

\textbf{Phần 2. Trắc nghiệm đúng sai}
\setcounterpageref{ex}{0}

\begin{ex} %[3D2V70]
Xét tính đúng sai của phép nhân số có 5 chữ số với số có 1 chữ số:
\choiceTF
{\True Nhân lần lượt từ phải sang trái (từ hàng đơn vị).}
{$20\,000 \times 5 = 10\,000$.}
{\True Tích của số có năm chữ số với 1 luôn bằng chính nó.}
{$15\,000 \times 4 = 50\,000$.}
\loigiai{
\begin{itemize}
    \item Ý a) Đúng: Quy tắc chuẩn.
    \item Ý b) Sai: Sai vì $20\,000 \times 5 = 100\,000$.
    \item Ý c) Đúng: Tính chất của số 1.
    \item Ý d) Sai: Sai vì $15\,000 \times 4 = 60\,000$.
\end{itemize}
\begin{itemize}
    \item a) \textbf{Đúng}.
    \item b) \textbf{Sai}.
    \item c) \textbf{Đúng}.
    \item d) \textbf{Sai}.
\end{itemize}
}
\end{ex}

\begin{ex} %[3D2V70]
Một xí nghiệp có 3 xưởng may, mỗi xưởng may được $14\,200$ bộ quần áo mỗi tháng:
\choiceTF
{\True Mỗi tháng xí nghiệp may được $42\,600$ bộ quần áo.}
{\True Trong 2 tháng một xưởng may được $28\,400$ bộ quần áo.}
{Cả xí nghiệp may trong 2 tháng được hơn 90 000 bộ.}
{\True Làm tròn sản lượng một tháng của xí nghiệp đến hàng nghìn là 43 000 bộ.}
\loigiai{
\begin{itemize}
    \item Ý a) Đúng: $14\,200 \times 3 = 42\,600$ bộ.
    \item Ý b) Đúng: $14\,200 \times 2 = 28\,400$ bộ.
    \item Ý c) Sai: Sai vì $42\,600 \times 2 = 85\,200 < 90\,000$ bộ.
    \item Ý d) Đúng: $42\,600$ làm tròn đến hàng nghìn là 43 000.
\end{itemize}
\begin{itemize}
    \item a) \textbf{Đúng}.
    \item b) \textbf{Đúng}.
    \item c) \textbf{Sai}.
    \item d) \textbf{Đúng}.
\end{itemize}
}
\end{ex}

\textbf{Phần 3. Trả lời ngắn}
\setcounterpageref{ex}{0}

\begin{ex} %[3D2V70]
Tính kết quả phép nhân: $15\,620 \times 5$.
\shortans{78100}
\loigiai{
\textbf{Đáp số:} $78100$\\
$15\,620 \times 5 = 78\,100$.
}
\end{ex}

\begin{ex} %[3D2V70]
Một hình vuông có cạnh dài $23\,500\text{ cm}$. Tính chu vi hình vuông đó theo xăng-ti-mét.
\shortans{94000}
\loigiai{
\textbf{Đáp số:} $94000$\\
$23\,500 \times 4 = 94\,000\text{ (cm)}$.
}
\end{ex}

\begin{ex} %[3D2V70]
Tìm số tự nhiên $y$, biết: $y : 3 = 24\,600$.
\shortans{73800}
\loigiai{
\textbf{Đáp số:} $73800$\\
$y = 24\,600 \times 3 = 73\,800$.
}
\end{ex}

\begin{ex} %[3D2V70]
Một nhà máy sản xuất mỗi ngày $16\,250$ sản phẩm. Hỏi 4 ngày nhà máy sản xuất bao nhiêu sản phẩm?
\shortans{65000}
\loigiai{
\textbf{Đáp số:} $65000$\\
$16\,250 \times 4 = 65\,000$ sản phẩm.
}
\end{ex}

\begin{ex} %[3D2V70]
Tính giá trị biểu thức: $12\,500 \times 6 - 25\,000$.
\shortans{50000}
\loigiai{
\textbf{Đáp số:} $50000$\\
$75\,000 - 25\,000 = 50\,000$.
}
\end{ex}

\begin{ex} %[3D2V70]
Có 4 kho lúa giống, mỗi kho chứa $18\,400\text{ kg}$. Hỏi 4 kho có tất cả bao nhiêu ki-lô-gam lúa giống?
\shortans{73600}
\loigiai{
\textbf{Đáp số:} $73600$\\
$18\,400 \times 4 = 73\,600\text{ (kg)}$.
}
\end{ex}
